\documentclass[12pt]{article}
\usepackage[utf8]{inputenc}
\usepackage[T1]{fontenc}
\usepackage{hyperref}
\title{Tsonga Language Resources for K-12 Curriculum Development}
\author{Compiled by Claude Sonnet 5 (claude-sonnet-5) for Andr\'as Kornai}
\date{}
\begin{document}
\maketitle

\section{Introduction}
Tsonga (Xitsonga; ISO 639-3 tso, Glottocode tson1249) is an official language of South Africa and is used as a language of instruction. This report rests on one search, and several generic ``guide'' PDFs it returned were unreliable and are not cited.

\section{Primary Educational Resources}
\begin{itemize}
\item Via Afrika Xitsonga Home Language Grade 5 Reader (CAPS-aligned; used alongside the Learner's Book and Teacher's Guide): \url{https://viaafrika.com/?p=10703}
\end{itemize}
The search also described a Department of Basic Education Xitsonga Home Language Grade 1 Book 1 (ISBN 978-1-920458-15-7), but no verified URL was returned for it, so it is not listed.

\section{Linguistic and Reference Materials}
\begin{itemize}
\item UNISA Xitsonga subject and module listing, including a grammar module (TSO3701) covering phonetics, phonology, morphology, syntax and semantics: \url{https://unisa.ac.za/sites/corporate/default/Register-to-study-through-Unisa/Subjects-&-modules/All-subjects/XITSONGA}
\end{itemize}
No Xitsonga dictionary or grammar book was found.

\section{Religious and Community Resources}
None were found in this search.

\section{Digital and Audio Resources}
None were found in this search.

\section{Curriculum Development Context}
Evidence of CAPS-aligned materials exists but is thin in these results. Dictionaries, DBE workbooks and other grades need further searching (PanSALB, university departments, DBE).
\end{document}